\section{Methods}
\label{sec:methods}
We compare random sampling, historical caching, fixed subsets, and adaptive
testing. Each defines an executed set $S_r^{(m)}$ and an estimator of the
complete pass rate.


\subsection{Random Sampling}
\label{sec:method-random}

Random sampling draws $k$ questions uniformly without replacement and reports
their unweighted pass rate:
\begin{equation}
  \hat{s}_r^{\mathrm{rand}}
  = \frac{1}{k}\sum_{q_i\in S_r^{\mathrm{rand}}}y_{ri}.
\end{equation}

\subsection{Historical Outcome Caching}
\label{sec:method-cache}
Questions that consistently pass or fail across recent runs are less likely to change
their outcomes after a modest agent update.
Historical caching leverages this stability by reusing expected verdicts for selected questions instead of executing them again.
From the available history $\mathcal{H}_r$, we
select a look-back window $\mathcal{W}_r$ and compute the historical
pass rate of each question $q_i$:
\begin{equation}
  \bar{y}_{ri}
  = \frac{1}{|\mathcal{W}_r|}
    \sum_{u\in\mathcal{W}_r}y_{ui}.
\end{equation}
A question is cache eligible when
\begin{equation}\label{equation:cache-eligible}
  \min(\bar{y}_{ri},1-\bar{y}_{ri}) \leq \epsilon,
\end{equation}
where $\epsilon\in[0,0.5]$ controls the required stability. Eligible questions
receive their recent majority verdict,
$\tilde{y}_{ri}=\mathbf{1}[\bar{y}_{ri}\geq 0.5]$, while all others are executed.
If $K_r$ is the cached set, the method executes
$S_r=\mathcal{Q}\setminus K_r$ and reports
\begin{equation}
  \hat{s}_r^{\mathrm{cache}}
  = \frac{1}{N}\left(
      \sum_{q_i\in S_r}y_{ri}
      + \sum_{q_i\in K_r}\tilde{y}_{ri}
    \right).
\end{equation}


\subsection{Fixed-Subset Evaluation}
\label{sec:method-fixed}

Inspired by \citet{maiapolo2024tinybenchmarks},
fixed-subset evaluation uses historical outcomes to select a representative set
that is reused across later runs. We compare three selection strategies and two
score estimators to separate subset construction from score reconstruction.

\subsubsection{IRT Models}
Several selectors and estimators use item response theory (IRT) models fitted to
the historical response matrix. A multidimensional two-parameter logistic (2PL)
model assigns each run an ability vector $\boldsymbol{\theta}_r$, each question
a discrimination vector $\boldsymbol{\alpha}_i$, and each question a scalar
difficulty $\beta_i$. The model predicts the probability of a pass as
\begin{equation}
  p_{ri}
  = \Pr(y_{ri}=1\mid\boldsymbol{\theta}_r,
                  \boldsymbol{\alpha}_i,\beta_i)
  = \sigma(\boldsymbol{\alpha}_i^\top\boldsymbol{\theta}_r-\beta_i),
\end{equation}
where $\sigma$ is the logistic sigmoid. We estimate parameters from historical
verdicts by maximum a posteriori estimation with Gaussian priors. The Rasch, or
one-parameter logistic (1PL), model is the unidimensional special case with
$\alpha_i=1$:
\begin{equation}\label{equation:rasch}
  p_{ri}=\sigma(\theta_r-\beta_i).
\end{equation}

\subsubsection{Subset Selection}
We use difficulty stratification as a direct, interpretable construction that selects questions across the estimated difficulty range. For comparison, we also evaluate
historical-response and IRT-feature clustering, adapting the two clustering selectors of
\citet{maiapolo2024tinybenchmarks}.


\paragraph{Difficulty stratification.}
We fit the Rasch model (Eq.\ref{equation:rasch}) on the calibration runs, sort questions by difficulty $\beta_i$,
and divide them into $k$ approximately equal strata. From each
stratum, we select the question closest to its median difficulty, and weight it by the stratum's fraction of the
benchmark.

\paragraph{Historical-response clustering.}
We represent each question by a vector recording its pass/fail outcomes across
historical runs and apply K-means clustering with $k$ clusters. The question
closest to each cluster centroid is selected, and its weight is the fraction of
benchmark questions in that cluster.

\paragraph{IRT-feature clustering.}
We instead represent each question by its fitted 2PL item parameters,
$[\boldsymbol{\alpha}_i;\beta_i]$, and apply the same K-means selection and weighting.

\subsubsection{Score Estimation}
Following \citet{maiapolo2024tinybenchmarks}, we evaluate each selector with two
score estimators. The weighted subset estimator uses only the selected
questions, while generalized p-IRT estimator (gp-IRT) reconstructs the unobserved questions
using the fitted IRT model.


\paragraph{Weighted subset estimator.}
This estimator treats each selected question as the representative of its
group.
For fixed subset $S$, let $w_i$ be the cluster or stratum weight of question
$q_i$, with $\sum_{q_i\in S}w_i=1$. The weighted estimate is
\begin{equation}
  \hat{s}_r^{\mathrm{weighted}}
  = \sum_{q_i\in S}w_i y_{ri}.
\end{equation}

\paragraph{Generalized p-IRT estimator.}

The gp-IRT estimator combines the weighted subset estimate with an IRT-based
reconstruction of the entire benchmark, balancing sampling variation against
error from the fitted IRT model.
The IRT-based reconstruction is p-IRT. It first estimates the current run’s ability
$\hat{\boldsymbol{\theta}}_r$ from its responses in $S$, holding fitted item
parameters fixed, and then replaces each unobserved verdict with its predicted pass
probability:
\begin{equation}\label{equation:pirt}
  \hat{s}_r^{\mathrm{pIRT}}
  = \frac{1}{N}\left(
      \sum_{q_i\in S}y_{ri}
      + \sum_{q_i\notin S}p_i(\hat{\boldsymbol{\theta}}_r)
    \right).
\end{equation}
In our experiments, difficulty stratification uses Rasch parameters for selection and
reconstruction, while the clustering variants use the fitted multidimensional 2PL
model.


The gp-IRT estimator blends the weighted and p-IRT estimates:
\begin{equation}
  \hat{s}_r^{\mathrm{gpIRT}}
  = \lambda\hat{s}_r^{\mathrm{weighted}}
    +(1-\lambda)\hat{s}_r^{\mathrm{pIRT}},
\end{equation}
where the blending weight 
$\lambda$ is:
\begin{equation}
  \lambda =
  \frac{\hat{b}^{2}}
       {\hat{\sigma}^{2}/(4|S|)+\hat{b}^{2}}.
\end{equation}
The weight favors the direct estimate when estimated IRT reconstruction bias
$\hat{b}^2$ is large relative to weighted-estimator variance
$\hat{\sigma}^2/(4|S|)$. We estimate $\hat{\sigma}^2$ by averaging within-run
question-outcome variance over historical runs and $\hat{b}^2$ as squared IRT
reconstruction MAE on held-out historical runs. Following the default
anchor-estimator adjustment of \citet{maiapolo2024tinybenchmarks}, we divide
$\hat{\sigma}^2$ by four, equivalently halving the estimated standard deviation.




\subsection{IRT-Based Adaptive Testing}
\label{sec:method-adaptive}
Adaptive testing selects each next question using responses observed in the
current run. The intuition is to always pick the questions
that are most informative for the current evaluation. We adapt the Rasch/Fisher-information procedure of
\citet{truong2025reliable} and additionally evaluate a multidimensional 2PL extension with
D-optimal selection.

We first calibrate item difficulties $\{\beta_i\}_{i=1}^N$ from historical runs
and initialize a new run’s ability at $\hat{\theta}_r^{(0)}=0$. At step $t$, the
predicted pass probability for each unobserved question is
\begin{equation}
  p_i^{(t)}=\sigma(\hat{\theta}_r^{(t)}-\beta_i).
\end{equation}
Under the Rasch model, the Fisher information supplied by question $q_i$ is
\begin{equation}
  \mathcal{I}_i(\hat{\theta}_r^{(t)})
  =p_i^{(t)}(1-p_i^{(t)}).
\end{equation}
At each step, the method executes the unobserved question with maximum
information and re-estimates $\hat{\theta}_r$ from all observed verdicts,
stopping after $k$ questions.

The multidimensional 2PL variant estimates ability vector
$\hat{\boldsymbol{\theta}}_r$ and uses D-optimality to reduce uncertainty across
its coordinates. Given the accumulated Fisher information matrix $\mathbf{I}_t$,
it selects the question maximizing
$p_i^{(t)}(1-p_i^{(t)})\boldsymbol{\alpha}_i^\top
\mathbf{I}_t^{-1}\boldsymbol{\alpha}_i$, then updates the ability vector and
information matrix from all observed verdicts.

Because Fisher selection targets informative rather than representative
questions, the raw pass rate of the selected questions need not approximate the
complete-benchmark pass rate. We therefore estimate
the complete pass rate using the p-IRT reconstruction defined above (Eq.\ref{equation:pirt})